\chapter{Reproducing the Results}

All commands run from the \texttt{haat-studio} folder. The report is rebuilt, with every number and figure read again from the result files, by
\begin{lstlisting}[language=bash,numbers=none]
ml/.venv/bin/python docs/mini-project-report/build.py --tests
\end{lstlisting}

\begin{table}[!htbp]
\centering
\caption{Where each result comes from.}
\label{tab:repro}
\small
\begin{tabularx}{\linewidth}{@{}L{42mm}Y@{}}
\toprule
\textbf{Result} & \textbf{Command (output file)}\\
\midrule
Generation times & made by the application; \texttt{data/results/*/meta.json}\\
Quota and run limit & made by the application; \texttt{data/gpu-quota.json}\\
E1--E3 & \texttt{ml/scripts/check\_real\_model.py} (\texttt{ml/runs/real-check/report.json})\\
E4 & \texttt{ml/scripts/lora\_microrun.py --scale 0.5 --steps 100} (\texttt{ml/runs/microrun/})\\
E5 & \texttt{ml/scripts/build\_pairs.py --retail <catalogue> --out ml/data/pairs} (\texttt{qc.json}, \texttt{pairs.csv})\\
Market analysis, engines & \texttt{analysis/market.py --retail <catalogue> --out analysis/results/market.json}\\
Tests & \texttt{npm test}; \texttt{cd ml \&\& .venv/bin/python -m pytest -q tests}\\
\bottomrule
\end{tabularx}
\end{table}

\paragraph{E6, on a CUDA GPU with at least 16~GB}
\begin{lstlisting}[language=bash,numbers=none]
cd ml
python scripts/build_pairs.py --retail <catalogue> --out data/pairs
python scripts/prepare_data.py --csv data/pairs/pairs.csv --out data/prepared
python scripts/train_lora.py --data data/prepared --out runs/indian-v1 --steps 2000
python scripts/evaluate.py --data data/prepared \
    --lora runs/indian-v1/lora-best.safetensors --out runs/indian-v1/eval
python scripts/compare_lora.py --run runs/indian-v1 \
    --out ../docs/mini-project-report/results/lora_comparison.json
cd .. && ml/.venv/bin/python docs/mini-project-report/build.py
\end{lstlisting}
\noindent The last command fills Section~\ref{sec:e6}, the abstract and the conclusions from \texttt{lora\_comparison.json}.

\chapter{Selected Source Code}

\section{The LoRA Layer}
From \texttt{ml/haat\_ft/lora.py}.
\begin{lstlisting}[language=python]
class LoRALinear(nn.Module):
    def __init__(self, base: nn.Linear, rank: int, alpha: float, dropout: float = 0.0):
        super().__init__()
        self.base = base
        self.base.requires_grad_(False)          # the original layer is frozen
        self.rank = rank
        self.scale = alpha / rank
        # float32 adapters over a bf16 base: tiny updates vanish in bf16
        self.down = nn.Linear(base.in_features, rank, bias=False, dtype=torch.float32)
        self.up = nn.Linear(rank, base.out_features, bias=False, dtype=torch.float32)
        nn.init.kaiming_uniform_(self.down.weight, a=math.sqrt(5))
        nn.init.zeros_(self.up.weight)           # B = 0: starts as the original model

    def forward(self, x):
        delta = self.up(self.down(x.to(self.down.weight.dtype))) * self.scale
        return self.base(x) + delta.to(x.dtype)
\end{lstlisting}

\section{The Reconstructed Training Loss}
From \texttt{ml/haat\_ft/flow.py}.
\begin{lstlisting}[language=python]
def sample_t(n, device, mu=1.5, std=1.0):
    return torch.sigmoid(torch.randn(n, device=device) * std - mu)

def flow_targets(image, t, noise=None):
    """Point on the noise-to-image line at time t, and the velocity along it."""
    noise = torch.randn_like(image) if noise is None else noise
    tt = t.view(-1, *([1] * (image.dim() - 1))).to(image.dtype)
    return (1 - tt) * noise + tt * image, image - noise

# inside flow_loss(): one coin per sample drops ALL conditions together
keep = torch.rand(b, device=target.device) >= p_uncond
pred = model(x_t, t, ca_images=batch["ca"], garment_images=batch["garment"],
             person_poses=batch["person_pose"], garment_poses=batch["garment_pose"],
             garment_categories=batch["category"], mask=keep)["x"]
w = batch["weight"] * (1 + (garment_weight - 1) * batch["garment_mask"])
err = (pred.float() - v.float()) ** 2 * w
loss = err.sum() / (w.sum() * target.shape[1]).clamp(min=1)
\end{lstlisting}

\section{Pose Guardrails for a Target Photograph}
From \texttt{ml/scripts/build\_pairs.py}.
\begin{lstlisting}[language=python]
def target_ok(p, category):
    if p["second_person"] > RULES["max_second_person"]:
        return False, "more than one person"
    if not (p["shoulders"] and p["hips"]):
        return False, "torso not visible"
    if category in ("bottoms", "one-pieces", "saree") and not p["knees"]:
        return False, "legs not visible"
    if p["frontal_ratio"] < RULES["min_shoulder_to_torso"]:
        return False, "not frontal"
    if not (p["face"] and p.get("facing_front")):
        return False, "back view or face cropped"
    if p["person_height"] < RULES["min_person_height"]:
        return False, "person too small"
    return True, ""
\end{lstlisting}

\section{The Quota Arithmetic}
From \texttt{src/lib/quota.ts}.
\begin{lstlisting}[language=java]
export const REQUEST_SECONDS = 60;   // what each call must have left to start

export function photosFitting(leftSeconds: number, perPhotoSeconds: number,
                              requestSeconds = REQUEST_SECONDS): number {
  return leftSeconds >= requestSeconds
    ? Math.floor((leftSeconds - requestSeconds) / perPhotoSeconds) + 1
    : 0;
}
// photosLeft = min(byTime, runsLimit - runsUsed); runsLimit is learned from
// Hugging Face's "runs limit" refusal.
\end{lstlisting}
